\documentclass[12pt, oneside, a4paper, brazil]{abntex2}
% \usepackage[utf8]{inputenc}
\usepackage[brazil]{babel}
% \usepackage[T1]{fontenc}
\usepackage{fontspec}
\usepackage{setspace}
\usepackage{graphicx}
\usepackage{scalefnt}
\usepackage{float}
\usepackage[a4paper, left=3cm, right=2.5cm, top=3cm, bottom = 2.5cm]{geometry}
\usepackage{hyperref}
\usepackage{listings}
\usepackage{color}
\usepackage{titlesec}
%\usepackage[alf]{abntex2cite}
% \usepackage{times}
\usepackage[scaled]{helvet}
\usepackage{fancyhdr}
\usepackage{lmodern}
\usepackage{lipsum}
\usepackage{abntex2cite}


%%% Definição de variáveis. %%%
% ---
% Informações de dados para CAPA e FOLHA DE ROSTO
% ---
\titulo{Análise de desempenho de estruturas de dados para um indexador de arquivos}
\autor{Atílio Antônio Dadalto \\ Tiago da Cruz Santos}
\local{Vitória}
\data{2019}
\instituicao{%
  Universidade Federal do Espírito Santo
  \par Departamento de Informática}
\tipotrabalho{Relatório}
\preambulo{Relatório apresentado como requisito parcial para aprovação na disciplina de Estrutura de Dados I, pela Universidade Federal do Espírito Santo.}

\newcommand{\versao}{2.0}
\newcommand{\subtitulo}{Anteprojeto}


%--- Configurações do environment lstlisting ---
\newfontfamily{\consolas}{Consolas.ttf}

\definecolor{gray}{rgb}{0.5,0.5,0.5}

\lstset{
  language=C,
  basicstyle=\footnotesize\consolas,
  numbers=left,
  numberstyle=\tiny,
  frame=tb,
  tabsize=4,
  breaklines=true,
  postbreak=\mbox{\textcolor{red}{$\hookrightarrow$}\space},
  columns=fixed,
  showstringspaces=false,
  showtabs=false,
  keepspaces,
  commentstyle=\color{gray},
  keywordstyle=\color{blue}
}

% Macros específicas do trabalho.
% (*) Inclua aqui termos que são utilizados muitas vezes e que demandam formatação especial.
% Os exemplos abaixo incluem i* (substituindo o asterisco por uma estrela) e Java com TM em superscript.
% Use sempre \xspace para que o LaTeX inclua espaço em branco após a macro somente quando necessário.
\newcommand{\istar}{\textit{i}$^\star$\xspace}
\newcommand{\java}{Java\texttrademark\xspace}
\newcommand{\latex}{\LaTeX\xspace}


%%% Configurações finais de aparência. %%%

% Altera o aspecto da cor azul.
\definecolor{blue}{RGB}{41,5,195}

% Informações do PDF.
\makeatletter
\hypersetup{
	pdftitle={\@title}, 
	pdfauthor={\@author},
	pdfsubject={\imprimirpreambulo},
	pdfcreator={LaTeX with abnTeX2},
	pdfkeywords={abnt}{latex}{abntex}{abntex2}{trabalho acadêmico}, 
	colorlinks=true,				% Colore os links (ao invés de usar caixas).
	linkcolor=blue,					% Cor dos links.
	citecolor=blue,					% Cor dos links na bibliografia.
	filecolor=magenta,				% Cor dos links de arquivo.
	urlcolor=blue,					% Cor das URLs.
	bookmarksdepth=4
}
\makeatother

% Espaçamentos entre linhas e parágrafos.
\setlength{\parindent}{1.3cm}
\setlength{\parskip}{0.2cm}


%%% Páginas iniciais do documento: capa, folha de rosto, ficha, resumo, tabelas, etc. %%%

% Compila o índice. <--- Desnecessário em Plano de Estudo.
%\makeindex
% Inicia o documento.
\begin{document}

% Retira espaço extra obsoleto entre as frases.
\frenchspacing


% Brasão da instituição.
\begin{figure}[h]
  \centering
  \includegraphics[scale=0.055]{figuras/brasao}
\end{figure} 

% Capa do trabalho.
\imprimircapa
\imprimirfolhaderosto

% Lista de silgas.
% (*) Indicar as siglas utilizadas no trabalho como no exemplo abaixo.
%\begin{siglas}
    %\item [UML] Linguagem de Modelagem Unificada, do inglês \textit{Unified Modeling Language}
%\end{siglas}

% Índice de capítulos.
\tableofcontents*


%%% Início da parte de conteúdo do documento. %%%

% Marca o início dos elementos textuais.
\clearpage
\textual

% Inclusão dos capítulos.
% (*) Para facilitar a organização, os capítulos foram divididos em arquivo separados e colocados dentro da.
% pasta capitulos/. Caso o aluno prefira trabalhar com um só arquivo, basta substituir os comandos \input 
% pelos conteúdos dos arquivos que estão sendo incluídos, excluindo a pasta capitulos/ em seguida.

\chapter*[INTRODUÇÃO]{INTRODUÇÃO}\label{cap-introducao} % Basicamente finalizado
\addcontentsline{toc}{chapter}{INTRODUÇÃO}

O presente documento tem como objetivo apresentar o relatório de resultados da implementação do nosso algoritmo para resolver o problema de agrupamento, no caso o agrupamento de espaçamento máximo, utilizando árvores geradoras mínimas (Minimum Spanning Tree ou MST).

No primeiro momento, apresentamos a organização das estruturas de dados escolhidas e o funcionamento do programa. Em seguida, uma breve análise de complexidade das principais etapas do processo. Por fim, a análise empírica para os cinco casos de teste disponibilizados.

Observamos nos resultados obtidos que o tempo de execução de cada etapa está em conformidade com a complexidade de cada passo. Para uma entrada pequena (por exemplo, a do caso 1), o tempo de leitura e escrita consomem quase $50\%$ do tempo total de execução do programa, visto que há considerável computação gasta em \emph{overhead} relacionado à manipulação de arquivos, enquanto a ordenação das distâncias consome cerca de $42\%$. É proveitoso observar, todavia, que quando a entrada cresce, o tempo de execução das etapas de leitura e escrita tendem a $0\%$ do tempo total de execução. Por outro lado, a etapa de ordenação das distâncias passa a ser cerca de $95\%$ do tempo total de execução, como podemos vislumbrar para os testes 3, 4 e 5 (vide Tabela~\ref{tab-tempo-execucao}). 
\chapter{Planejamento e implementação}\label{cap-planejamento-implementacao}

O planejamento das soluções foi feito tendo em mente o objetivo final de um problema. As soluções foram subsequentemente divididas em passos cada vez menores e mais discretizados, possibilitando a abordagem de dividir para conquistar.

A implementação, por sua vez, conta com vários usos de expressões lambda e também da função \texttt{map}. O paradigma recursivo foi usado em situações que julgou-se mais simples ou intuitivo sua implementação, como por exemplo na ordenação ou na remoção de elementos duplicados em uma lista (mais sobre isso nas Seções~\ref{problemaB} e \ref{problemaC}). Para maior entendimento das soluções, é claro, é necessário visualizar os códigos, que aqui foram omitidos. Todo o projeto busca seguir as convenções da PEP-8. 

\section{Problema A}\label{problemaA}
\begin{enumerate}[label=\textbf{\alph*)},series=problemas]
\item Calcular a distância percorrida por um determinado robô ao longo do processo de resgate das vítimas. Considere que a distância total percorrida deve ser calculada como a soma de todas as distâncias entre os pontos de passagem do robô.
\end{enumerate}

\noindent \textbf{COMPREENSÃO DO PROBLEMA E PLANEJAMENTO}: deve-se extrair a tupla de um robô da lista de entrada, de acordo com o identificador dado. Em seguida, calcular a distância total que o robô percorreu e retornar esse valor.

Sendo assim, podemos criar as seguintes funções para solucionar o problema:
\begin{enumerate}
    \item Função para cálculo de distância no plano cartesiano
    \item Função para extrair todas as tuplas de um mesmo robô da lista
    \item Função para extrair todos os pontos percorridos por um mesmo robô 
    \item Função para calcular a distância total de uma lista de pontos
\end{enumerate}

A primeira, \texttt{dist\_euclid}, é trivial e foi implementada diversas vezes durante o curso;

Na segunda, \texttt{lista\_tuplas\_robo\_id}, foi utilizada compreensão de lista com condicional para formar uma lista de tuplas;

Na terceira, \texttt{pega\_pontos\_robo}, foi aplicada, a todos os elementos da lista de robôs, a função que retorna o local do robô dada uma tupla de robô, por meio da função \texttt{map}. A partir daí, adiciona-se a origem à lista de pontos percorridos, pois todos os robôs partem da origem.

Por fim, a quarta, \texttt{distancia\_total}, é utilizada a recursão para calcular a distância de um ponto até outro, partindo do primeiro ponto, dada uma lista de pontos como entrada.

Com essas quatro funções, temos o suficiente para solucionar o Problema A; a função que propriamente resolve essa questão é a \texttt{distancia\_total\_robo}, ao extrair todos os pontos de um robô dado seu id e a lista de robôs e calcular a distância do primeiro ponto até o último, passando por todos os pontos no meio.


\section{Problema B}\label{problemaB}
\begin{enumerate}[label=\textbf{\alph*)},resume*=problemas]
\item Determine qual dos robôs apresenta o seu último ponto de passagem no terreno de busca que possui a maior distância em relação à origem. Exiba o caminho percorrido pelo robô e o tempo total do percurso;
\end{enumerate}

\noindent \textbf{COMPREENSÃO DO PROBLEMA E PLANEJAMENTO}: deve-se extrair, da lista de entrada, o último ponto de passagem de todos os robôs. Com isso, determinar qual robô possui ponto mais longe da origem e então imprimir caminho percorrido pelo robô, além de determinar tempo do percurso.

Sendo assim, podemos criar as seguintes funções para solucionar o problema:
\begin{enumerate}
    \item Função para obter uma lista com os últimos pontos de todos os robôs da lista de entrada 
    \item Função para calcular a distância da origem a todos esses pontos, retornando uma lista 
    \item Função para buscar os robôs que obtiveram a maior distância
    \item Função para determinar tempo de percurso de um robô
\end{enumerate}

A primeira função, \texttt{ultimos\_pontos\_robos}, é muito simples, bastando aplicar a função \texttt{map} com uma função que retorne o último elemento de uma lista e a lista de tuplas de robôs.

A segunda, \texttt{dist\_origem\_ultimo\_ponto}, também é simples e utiliza a função de cálculo de distância da origem até um ponto com todos os pontos obtidos pela primeira função, por meio da função \texttt{map}.

A terceira, \texttt{indices\_ids\_mais\_distantes}, após calcular as distâncias até os últimos pontos de todos os robôs, calcula os índices dos robôs que tiveram maior distância (caso exista mais de um). Esses índices podem ser utilizados para acessar a lista de ids que a função \texttt{ids\_robos} retorna, função essa que retorna uma lista com todos os ids dos robôs, sem repetição.

Finalmente, a quarta, \texttt{tempo\_percurso}, simplesmente retorna o instante da última ocorrência de um robô, já que isso representa o tempo total de percurso atual de um robô.

Sendo assim, teremos os ids dos robôs mais distantes, os percursos destes --- através da função \texttt{pega\_pontos\_robo} do Problema A com os ids supracitados --- e o tempo de percurso.


\section{Problema C}\label{problemaC}

\begin{enumerate}[label=\textbf{\alph*)},resume*=problemas]
\item Exiba os caminhos percorridos por todos os robôs que entraram no terreno de busca, ordenados crescentemente pela distância total percorrida;
\end{enumerate}

\noindent \textbf{COMPREENSÃO DO PROBLEMA E PLANEJAMENTO}: deve-se calcular as distâncias percorridas por todos os robôs, adicionando-as em uma tupla que será ordenada a partir das distâncias, impressa na tela e então retornada.

Sendo assim, podemos criar as seguintes funções para solucionar o problema:
\begin{enumerate}
\item Função para juntar id, distância total e caminho percorrido de todos os robôs em uma lista de tuplas
\item Função para ordenar essa lista pelas distâncias nas tuplas
\end{enumerate}

A primeira função simplesmente percorre as listas de ids, distâncias totais e caminhos, juntando elemento a elemento em uma tupla e concatenando em uma lista, de forma recursiva. A lista de distâncias totais é calculada utilizando a função \texttt{distancia\_total} do Problema A nas listas de pontos obtidas a partir da lista de robôs. Nas seções anteriores, discorre-se sobre a obtenção da lista de caminhos percorridos pelos robôs e da lista de IDs.

Para a segunda, que é a principal função do problema, foi utilizado o método de ordenação \textit{merge sort}, em sua variação recursiva. Além disso, o \textit{merge sort} implementado recebe uma lista de tuplas, das quais os segundos elementos são utilizados como referência para a ordenação da lista. O segundo elemento da tupla, no caso deste problema, será a distância total do robô em questão.

Posto isto, basta chamar a função de ordenação com a saída da primeira função e o Problema C estará solucionado.

\section{Problema D}\label{problemaD}

\begin{enumerate}[label=\textbf{\alph*)},resume*=problemas]
\item Forneça a identidade do(s) robô(s) que conseguiu(ram) informar o maior número de vítimas (considerando que não há duplicação de identificação de vítima por um mesmo robô).
\end{enumerate}

\noindent \textbf{COMPREENSÃO DO PROBLEMA E PLANEJAMENTO}: deve-se calcular o número de vítimas que cada robô avistou e identificar o robô que avistou o maior número.
Retornar mais de um robô caso o maior número seja repetido.

Sendo assim, podemos criar as seguintes funções para solucionar o problema:
\begin{enumerate}
    \item Função para extrair o número de vítimas vistas em cada ponto passado por um robô
    \item Função para somar elementos de uma lista  
    \item Função que aplique essas duas últimas funções para todos os robôs da lista de entrada  
    \item Função para extrair o maior número da lista obtida pelo último item  
    \item Função para obter IDs dos robôs que obtiveram o máximo calculado anteriormente
\end{enumerate}

Para a primeira função, \texttt{lista\_vitimas\_robo}, podemos lançar mão da compreensão de lista mais uma vez para, em cada ocorrência (elemento) da lista de robôs, extrairmos o número de vítimas avistadas nessa ocorrência se o robô da tupla em questão tiver o id de entrada.

A segunda, \texttt{soma\_lista}, é trivial e foi implementada diversas vezes durante o curso;

A terceira, \texttt{total\_vitimas\_robo}, como já mencionado, apenas utiliza a primeira função para obter a lista com vítimas avistadas por um robô e, com a segunda função, soma a lista para calcular o número total de vítimas avistadas por um robô.

A quarta, \texttt{max\_lista}, utiliza a função \texttt{reduce} com a função \texttt{maior} (que compara dois elementos e retorna o maior) e a lista de entrada para comparar elemento por elemento na lista, retornando o maior.

Agora, com a terceira, podemos utilizar a função \texttt{map} para extrair todas as vítimas de cada robô da lista de robôs, com auxílio da função \texttt{ids\_robos} e uma função anônima. Teremos então uma lista com o total de vítimas de cada robô. Utilizando a função \texttt{indices\_maximos} do Problema B, que retornará os índices de todos os elementos que coincidam com o maior elemento, podemos acessar a lista de ids com esses índices e retornar os ids dos robôs que avistaram o maior número de vítimas (se houver mais de um), solucionando, portanto, o Problema D.



\chapter*{CONCLUSÃO}\label{cap-conclusao}
\addcontentsline{toc}{chapter}{CONCLUSÃO}

% Discutir as conclusões e limitações do seu trabalho.

Pelo estudo realizado neste trabalho, fica evidente como podemos construir diversas soluções tendo como ponto de partida o paradigma funcional. Utilizando a modularização e o reuso de funções através de importações, foi possível conceber uma central de processamento fictícia adaptada para o contexto proposto ao projeto. Ademais, a compreensão de problemas e planejamento de soluções permitiram que o desenvolvimento das resoluções fosse muito mais focado e centrado. 

Os conceitos de paradigma recursivo e aplicativo, além da compreensão de lista foram fundamentais para a execução das soluções idealizadas. A utilização dessas abstrações permitiu uma solução mais limpa e sucinta do que se fosse empregado o paradigma procedural, por exemplo. O uso de recursão, no entanto, pode causar \textit{stack overflow} e, embora a implementação do \textit{merge sort} seja significativamente mais simples neste paradigma, um grande volume de dados poderia causar a suspensão da execução do programa; o próprio Python limita o nível de recursões a 1000 chamadas empilhadas por padrão, embora seja possível aumentar esse número.
\chapter{Análise  e resultados}\label{cap-analise-resultado}

Os testes aqui mostrados se referem aos resultados obtidos pela Máquina 1.\footnote{Planilha disponível no \href{https://docs.google.com/spreadsheets/d/1yzVdhAuRN6BSb2E2VDG79fv59cQ-DlLtYV5bvRnVaVQ/edit?usp=sharing}{Google Planilhas}} Os testes foram realizados rodando o indexador com uma estrutura por vez, e, como supracitado, realizando 50 iterações do mesmo teste. Os arquivos utilizados são 10x maiores a cada teste para que os gráficos possam ser analisados de forma linear, posto que o eixo y possui escala logarítmica. Com arquivos de 1GB, a Máquina 1 começa a enfrentar problemas para carregar estruturas, precisando lançar mão da memória swap --- substancialmente mais lenta que a RAM --- sobretudo com a árvore de prefixos, uma vez que esta estrutura demanda mais memória que as demais, em geral.

\section{Inserções}

\begin{center}\label{graf_ins_curva}
    \includegraphics[scale=0.7]{Relatório/figuras/graficos_m1/insercoes_curva.pdf}
\end{center}

\begin{center}\label{graf_ins_barra}
    \includegraphics[scale=0.7]{Relatório/figuras/graficos_m1/insercoes_barra.pdf}
\end{center}

\begin{center}\label{graf_ins_barra2}
    \includegraphics[scale=0.7]{Relatório/figuras/graficos_m1/insercoes_barra2.pdf}
    \small{Fonte: elaborado pelos autores.}
\end{center}

A discrepância entre os tempos da lista em comparação com o resto das estruturas talvez seja melhor visualizado com o eixo y em escala linear:
\begin{center}\label{graf_ins_eixoylinear}
    \includegraphics[scale=0.7]{Relatório/figuras/graficos_m1/insercoes_eixoy_linear.pdf}
\end{center}

Como visto acima, a lista encadeada chega a apresentar desempenho quase 50x pior se comparada à outras estruturas para leitura de arquivos de 10mb, mesmo tendo um desempenho equivalente às outras para arquivos de tamanho $\leq$ 1kb.

A tabela hash, por apresentar um tamanho $M$ expressivo, deve alocar muito espaço inicialmente, algo que diminui sua velocidade consideravelmente com arquivos pequenos. Para arquivos maiores, no entanto, a estrutura empatou com a árvore trie, cujo desempenho superou os das outras árvores.

\section{Buscas}\label{sec-grafbuscas}

\begin{center}\label{graf_buscas_curva}
    \includegraphics[scale=0.7]{Relatório/figuras/graficos_m1/buscas_curva.pdf}
\end{center}

\begin{center}\label{graf_buscas_barra}
    \includegraphics[scale=0.7]{Relatório/figuras/graficos_m1/buscas_barra.pdf}
\end{center}

\begin{center}\label{graf_buscas_barra2}
    \includegraphics[scale=0.7]{Relatório/figuras/graficos_m1/buscas_barra2.pdf}
    \small{Fonte: elaborado pelos autores.}
\end{center}

Para buscas, tivemos um empate entre a árvore binária não balanceada e a AVL. A tabela hash e a árvore de prefixos também obtiveram resultados semelhantes, com uma pequena margem de vitória para a árvore trie. Esta última atingiu os melhores tempos de buscas dentre todas as estruturas, com perda apenas em arquivos minúsculos (100b ou menos).
\chapter*{CONCLUSÃO}\label{cap-conclusao}
\addcontentsline{toc}{chapter}{CONCLUSÃO}

Com os dados levantados e a análise feita, podemos então discutir os resultados de uma forma geral. A amostra de 50 iterações é boa o bastante para tomarmos conclusões sobre os desempenhos das estruturas.

% por conta de sua inserção que percorre a lista toda em busca de um elemento que seja igual ao elemento a ser inserido 
% Mesmo que as inserções das árvores possuam complexidade computacional grande, estas possuem uma buscas mais velozes que a da lista, por exemplo.
Sendo assim, a lista encadeada, como mostrado no Capítulo~\ref{cap-analise-resultado}, é a estrutura mais lenta dentre todas do projeto, possuindo custo $O(n)$ em vez do $O(1)$ usual para inserções de listas encadeadas. Apesar de tudo, a implementação da lista também é extremamente mais simples comparada às implementações das árvores, o que não ajuda apenas a construção inicial de um projeto, mas também sua manutenção.

A árvore binária sem balanceamento teve desempenho quase idêntico à árvore AVL, vencendo por pouco no tempo de inserções e empatando no tempo de buscas.

A árvore binária balanceada, por sua vez, demonstrou um desempenho equilibrado dentre as 5 estruturas do projeto, tendo ainda uma implementação não extremamente complexa; por conta de seu equilíbrio, foi escolhida para acomodar Palavras na tabela hash.

A árvore de prefixos, ou trie, teve, em média, o melhor desempenho de todos para inserções e buscas e praticamente empatando com a tabela hash. No entanto, ela demanda muita memória, mesmo cada nó possuindo apenas 36 filhos.

Por fim, a tabela hash é uma estrutura peculiar, visto que o seu tamanho M pode variar enormemente de uma implementação para outra. O valor escolhido para nosso indexador, 100003, fez com que a tabela hash tivesse um desempenho inferior a outras estruturas para entradas pequenas, porém tendo certo ganho com aumento dos arquivos. Cada caso de uso da tabela hash deve ser analisado individualmente para que seja possível encontrar um número mais próximo do ideal, economizando, assim, tempo de execução e memória.

\enlargethispage{\baselineskip}
\enlargethispage{\baselineskip}
Portanto, cabe ao idealizador de um projeto determinar qual estrutura é mais adequada ao problema em questão, analisando, se possível, a razão entre inserção e busca e também se a solução não é complexa demais para o objetivo proposto. Para um projeto diminuto, a lista encadeada poderia ser utilizada sem problemas notáveis de performance, ainda sendo válido destacar sua simples implementação; no caso de um indexador, no entanto, ela representou a pior estrutura. Como visto na Seção~\ref{sec-grafbuscas}, as melhores estruturas para inserções e buscas, neste caso, foram a árvore trie e a tabela hash. Para um projeto como este de um indexador de palavras, os autores consideram o TAD árvore de prefixos como o melhor, estando tecnicamente empatada com a tabela hash.
\begin{apendicesenv}

\chapter{Planos de testes}\label{apendiceA}

Este apêndice serve como repositório para os planos testes das funções mais importantes implementadas no projeto. Algumas funções assumem certos valores de entrada, de forma a afunilar as possíveis entradas, dado que os valores a serem recebidos podem variar enormemente e que o projeto conta com muitas funções; no entanto, nos contemos a assumir entradas de maneira muito razoável, e, além disso, todos os casos básicos foram tratados devidamente dentro das funções.


Para melhor exemplificação e legibilidade dos testes das funções abaixo, usaremos a lista de dados de entrada ilustrada a seguir:

\begin{lstlisting}
listaRobos = [
    ('robo3', 1, (7, 7), 3), ('robo4', 2, (7, 5), 2),
    ('robo3', 3, (5, 4), 3), ('robo3', 4, (8, 1), 4),
    ('robo4', 5, (4, 5), 3), ('robo5', 6, (7, 7), 4)
]
\end{lstlisting}

\section{\texttt{distancia\_total()}}
A função calcula a distância total entre uma lista de pontos somando a distância entre um ponto e o próximo até o fim da lista.

\begin{center}
    \begin{tabular}{|c|c|}
        \cline{1-2}
        lista & resultado esperado \\ \hline
        [] & None \\ \hline  
        lista com um ponto & None \\ \hline  
        [p1, p2, p3] & distância de p1 até p2 + distância de p2 até p3 \\ \hline
    \end{tabular}
\end{center}

Definindo entradas:

\begin{center}
    \begin{tabular}{|c|c|c|}
        \cline{1-3}
        lista & resultado esperado & resultado obtido  \\ \hline
        [] & None & None \\ \hline
        [(4,4)] & None & None \\ \hline
        [(0,0),(3,4),(5,6)] & 7.82842712474619 & 7.82842712474619 \\ \hline

    \end{tabular}
\end{center}

\section{\texttt{distancia\_total\_robo()}}
A função calcula a distância total percorrida por um robô, ou seja, a soma das distâncias de todos os pontos percorridos por ele.

\begin{center}
    \begin{tabular}{|c|c|}
        \cline{1-2}
        lista & resultado esperado \\ \hline
        [] & None \\ \hline  
        [(a,(p1)),(a,(p2))], 'a' & distância de (0,0) até p1 + distância de p1 até p2 \\ \hline
    \end{tabular}
\end{center}


Definindo entradas:

\begin{center}
    \begin{tabular}{|c|c|c|}
        \cline{1-3}
        lista & resultado esperado & resultado obtido  \\ \hline
        [] & None & None \\ \hline
        [('robo1',2,(4,4),5)], 'robo2' & None & None \\ \hline
        listaRobos, 'robo3' & 17.74768689919494 & 17.74768689919494 \\ \hline

    \end{tabular}
\end{center} 


\section{\texttt{indices\_maximos()}}
A função recebe uma lista numérica e retorna uma lista com os índices de todas as ocorrências do valor máximo na lista.

\begin{center}
    \begin{tabular}{|c|c|}
        \cline{1-2}
        lista & resultado esperado \\ \hline
        [] & [] \\ \hline  
        [a, b, c, d] & índices das ocorrências do maior elemento da lista  \\ \hline
    \end{tabular}
\end{center}

Definindo entradas:

\begin{center}
    \begin{tabular}{|c|c|c|}
        \cline{1-3}
        lista & resultado esperado & resultado obtido  \\ \hline
        [] & [] & [] \\ \hline
        [5] & [0] & [0] \\ \hline
        [1,2,3,10,1,10] & [3,5] & [3,5] \\ \hline

    \end{tabular}
\end{center} 


\section{\texttt{imprime\_robos\_mais\_distantes()}}
A função recebe uma lista de robôs e imprime os robôs que encontram-se mais distantes da origem, seus percursos e tempos de percurso
\begin{center}
    \begin{tabular}{|c|c|}
        \cline{1-2}
        lista & resultado esperado \\ \hline
        [] & None \\ \hline  
        [('b',(10,10),1)] & ['b'],[(0,0),(10,10)],[1] \\ \hline
        [('a',(5,4),1),('b',(10,10),1)] & ['b'],[(0,0),(10,10)],[1]  \\ \hline
    \end{tabular}
\end{center}

Definindo entradas:

\begin{center}
    \begin{tabular}{|c|c|c|}
        \cline{1-3}
        lista & resultado esperado & resultado obtido  \\ \hline
        [] & None & None \\ \hline
        [('robo',1,(10,10),3)] & [('robo',((0,0),(10,10)),1] & [('robo',((0,0),(10,10)),1] \\ \hline
        [('robo',1,(10,10),3),('robo2',1,(2,1),3)] & [('robo',((0,0),(10,10)),3] & [('robo',((0,0),(10,10)),3]  \\ \hline

    \end{tabular}
\end{center}




\section{\texttt{caminhos\_robos\_crescente()}}
A função retorna uma lista de tuplas para cada robô, ordenadas em ordem crescente pela distância total percorrida pelos robôs.

\begin{center}
    \begin{tabular}{|c|c|}
        \cline{1-2}
        Entradas & resultado esperado \\ \hline
        [] &  None\\ \hline
        [(a),(b)] & soma do caminho percorrido a, soma do caminho b  \\ \hline
    \end{tabular}
\end{center}


Definindo entradas (saídas arredondadas para melhor visualização):

\begin{center}
    \begin{tabular}{|c|c|}
        \cline{1-2}
        lista & resultado esperado \\ \hline
        [] & None \\ \hline
        [('robo1',2,(4,4),5)] & [('robo1', 5.6, [(0, 0), (4, 4)])] \\ \hline
        [('robo1',2,(4,4),5),('robo3',7,(5,5),2)] & [('robo1', 5.6, [(0, 0), (4, 4)]), ('robo3', 7, [(0, 0), (5, 5)])] \\ \hline
    \end{tabular}
\end{center} 

\begin{center}
    \begin{tabular}{|c|}
        \cline{1-1}
        resultado obtido \\ \hline
        None \\ \hline
        [('robo1', 5.6, [(0, 0), (4, 4)])] \\ \hline
        [('robo1', 5.6, [(0, 0), (4, 4)]), ('robo3', 7, [(0, 0), (5, 5)])] \\ \hline
    \end{tabular}
\end{center}



\section{\texttt{merge\_ordenada\_tupla()}}
Dadas duas listas ordenadas pelo segundo elemento da tupla, junta-as de forma ordenada.

\begin{center}
    \begin{tabular}{|c|c|}
        \cline{1-2}
        lista & resultado esperado \\ \hline
        [] & [] \\ \hline
        [],[(a,b)] & [(a,b)] \\ \hline
        [(c,d)],[] & [(c,d)] \\ \hline
        [(a,b)],[(c,d)]  & [(a,b),(c,d)] \\ \hline
        [(a,b),(c,d)],[(b,c),(d,e)] & [(a,b),(b,c)(c,d),(d,e)] \\ \hline
    \end{tabular}
\end{center}

Definindo entradas:

\begin{center}
    \begin{tabular}{|c|c|c|}
        \cline{1-3}
        lista & resultado esperado & resultado obtido  \\ \hline
        [],[] & [] & []\\ \hline
        [(1,2)],[(3,4)] & [(1,2),(3,4)] & [(1,2),(3,4)] \\ \hline
        [(1,2),(3,4)],[(5,6),(7,8)] & [(1,2),(3,4),(5,6),(7,8)] & [(1,2),(3,4)],[(5,6),(7,8)] \\ \hline
    \end{tabular}
\end{center}


\section{\texttt{merge\_sort\_tupla()}}
Ordena uma lista de tuplas através do método \textit{merge sort}, tendo como referência o segundo elemento das tuplas.

\begin{center}
    \begin{tabular}{|c|c|}
        \cline{1-2}
        lista & resultado esperado \\ \hline
        [] & [] \\ \hline
        [(a)]  & [(a)] \\ \hline
        [(a,d),(c,b),(b,f)] & [(c,b),(a,d),(b,f)]  \\ \hline
    \end{tabular}
\end{center}


Definindo entradas:

\begin{center}
    \begin{tabular}{|c|c|c|}
        \cline{1-3}
        lista & resultado esperado & resultado obtido  \\ \hline
        [] & [] & [] \\ \hline
        [(1,2)] & [(1,2)] & [(1,2)] \\ \hline
        [(1,4),(2,3)] & [(2,3),(1,4)] & [(2,3),(1,4)] \\ \hline

    \end{tabular}
\end{center}


\section{\texttt{ids\_mais\_vitimas()}}
 A função retorna o(s) id(s) do(s) robô(s) que avistaram o maior número de vítimas.
 
 \begin{center}
    \begin{tabular}{|c|c|}
        \cline{1-2}
        lista & resultado esperado \\ \hline
        [] & None \\ \hline
        um robô  & o id do único robô \\ \hline
        mais de um robô, um robô com mais vítimas avistadas & id do robô com mais vítimas  \\ \hline
        mais de um robô, mais de um robô com mais vítimas & ids dos robôs com mais vítimas \\ \hline
    \end{tabular}
\end{center}

Definindo entradas:

\begin{center}
    \begin{tabular}{|c|c|c|}
        \cline{1-3}
        lista & resultado esperado & resultado obtido  \\ \hline
        [] & None & None \\ \hline
        [('robo1',1,(1,1),2)] & ['robo1'] & ['robo1'] \\ \hline
        [('robo1',1,(1,1),2),('robo2',5,(3,2),5)] & ['robo2'] & ['robo2'] \\ \hline
        [('robo1',1,(1,1),2),('robo2',5,(3,2),2)] & ['robo1', 'robo2'] & ['robo1', 'robo2'] \\ \hline
    \end{tabular}
\end{center} 

\end{apendicesenv}



% Finaliza a parte no bookmark do PDF para que se inicie o bookmark na raiz e adiciona espaço de parte no sumário.
\phantompart

% Marca o fim dos elementos textuais.
\postextual

% Referências bibliográficas
\clearpage
\bibliography{bibliografia}

% Fim do documento.
\end{document}
